% =====================================================================
% 06_security.tex
% =====================================================================
\begingroup
\newcolumntype{L}[1]{>{\hsize=#1\hsize\raggedright\arraybackslash}X}

\section{Bảo mật mạng WLAN}
\label{sec:security}

\subsection{Vì sao mạng không dây cần bảo mật mạnh?}

Bất kỳ ai nằm trong vùng phủ sóng vô tuyến đều có thể thu bắt các khung dữ
liệu. Vì vậy, một \ac{wlan} an toàn cần bảo đảm \emph{tính bí mật} (mã hóa),
\emph{tính toàn vẹn} (phát hiện sửa đổi), \emph{xác thực} (ai đang kết nối và
kết nối tới đâu) và \emph{tính sẵn sàng} (khả năng chống gây nhiễu và tấn công
khung quản lý).

\subsection{Các thế hệ bảo mật Wi-Fi}

\begin{table}[htbp]
  \centering\footnotesize
  \caption{Các thế hệ cơ chế bảo mật WLAN.}
  \label{tab:security-evolution}
  \begin{tabularx}{\linewidth}{@{} L{0.5} L{0.3} L{1.0} L{2.2} @{}}
    \toprule
    \textbf{Cơ chế} & \textbf{Năm} & \textbf{Mã mật / quản lý khóa} & \textbf{Tình trạng} \\
    \midrule
    \acs{wep} & 1997 & RC4, IV 24 bit, khóa tĩnh & Đã bị phá; có thể khôi phục khóa trong vài phút \cite{fms2001}; không sử dụng. \\
    \acs{wpa} & 2003 & RC4 với \acs{tkip}, \acs{psk} hoặc 802.1X & Giải pháp chuyển tiếp đã lỗi thời. \\
    WPA2 & 2004 & \acs{aes}-\acs{ccmp}, \acs{psk} hoặc 802.1X & Được triển khai rộng; mật khẩu yếu có thể bị dò ngoại tuyến và còn lỗi dạng \acs{krack}. \\
    WPA3 & 2018 & \acs{aes}-\acs{ccmp}/\acs{gcmp}, \acs{sae} hoặc 802.1X, bắt buộc \acs{pmf} & Khuyến nghị hiện hành \cite{wpa3spec}. \\
    \bottomrule
  \end{tabularx}
\end{table}

\subsection{WPA2 và quá trình bắt tay bốn bước}

Trong WPA2-Personal, cụm mật khẩu (8--63 ký tự) và \ac{ssid} được đưa vào
PBKDF2 để tạo khóa \ac{pmk} dài 256 bit:
\begin{equation}
  \mathrm{PMK} = \mathrm{PBKDF2}\bigl( \text{HMAC-SHA1},\ \text{passphrase},\ \text{SSID},\ 4096,\ 256 \bigr).
\end{equation}
Trong WPA2/WPA3-Enterprise, \ac{pmk} được tạo từ quá trình xác thực
802.1X/\ac{eap}. Quá trình bắt tay bốn bước (Hình~\ref{fig:4way}) xác nhận cả
hai phía đều biết \ac{pmk}, đồng thời dẫn xuất khóa phiên mới \ac{ptk} từ
\ac{pmk}, hai địa chỉ MAC và hai số ngẫu nhiên (nonce):
\begin{equation}
  \mathrm{PTK} = \mathrm{PRF}\bigl( \mathrm{PMK},\ \text{``Pairwise key expansion''},\ \mathrm{AA}\,\|\,\mathrm{SPA}\,\|\,\mathrm{ANonce}\,\|\,\mathrm{SNonce} \bigr),
\end{equation}
Trong đó, AA và SPA lần lượt là địa chỉ MAC của bộ xác thực và thiết bị cần
xác thực; chúng được sắp theo thứ tự min/max cố định để hai phía tính ra cùng
một giá trị. \ac{ptk} được chia thành khóa xác nhận, khóa mã hóa khóa và khóa
tạm thời dùng bởi \ac{ccmp}. Thông điệp~3 cũng gửi \ac{gtk} dùng cho lưu lượng
quảng bá/đa hướng; khóa này được mã hóa và bảo vệ bằng \ac{mic}.

\begin{figure}[htbp]
  \centering
  \begin{tikzpicture}[font=\small,>=Stealth,node distance=0.8cm]
    \node[font=\bfseries] (sta) {STA (thiết bị cần xác thực)};
    \node[font=\bfseries, right=6.0cm of sta] (ap) {AP (bộ xác thực)};

    \node[draw, fill=yellow!20, font=\scriptsize, below=0.4cm of sta] (pmk1)
    {Đã có PMK};
    \node[draw, fill=yellow!20, font=\scriptsize, below=0.4cm of ap] (pmk2)
    {Đã có PMK};

    \coordinate[below=6.5cm of pmk1] (sta_bottom);
    \coordinate[below=6.5cm of pmk2] (ap_bottom);
    \draw[thick] (pmk1.north) -- (sta_bottom);
    \draw[thick] (pmk2.north) -- (ap_bottom);

    % M1
    \coordinate[below=0.6cm of pmk1] (m1_l);
    \coordinate[below=0.6cm of pmk2] (m1_r);
    \draw[->] (m1_r) -- (m1_l) node[midway, above,
    font=\scriptsize]{M1: ANonce};
    \node[anchor=west, font=\scriptsize, below=0.1cm of m1_l] {dẫn xuất PTK};

    % M2
    \coordinate[below=1.4cm of m1_l] (m2_l);
    \coordinate[below=1.4cm of m1_r] (m2_r);
    \draw[->] (m2_l) -- (m2_r) node[midway, above,
    font=\scriptsize]{M2: SNonce, MIC, thông số bảo mật};
    \node[anchor=east, font=\scriptsize, below=0.1cm of m2_r] {dẫn xuất PTK,
    kiểm tra MIC};

    % M3
    \coordinate[below=1.4cm of m2_l] (m3_l);
    \coordinate[below=1.4cm of m2_r] (m3_r);
    \draw[->] (m3_r) -- (m3_l) node[midway, above,
    font=\scriptsize]{M3: GTK (đã mã hóa), MIC, yêu cầu cài đặt khóa};
    \node[anchor=west, font=\scriptsize, below=0.1cm of m3_l] {kiểm tra MIC,
    cài đặt PTK/GTK};

    % M4
    \coordinate[below=1.4cm of m3_l] (m4_l);
    \coordinate[below=1.4cm of m3_r] (m4_r);
    \draw[->] (m4_l) -- (m4_r) node[midway, above,
    font=\scriptsize]{M4: ACK, MIC};
    \node[anchor=east, font=\scriptsize, below=0.1cm of m4_r] {cài đặt PTK};
  \end{tikzpicture}
  \caption{Quá trình bắt tay bốn bước trong WPA2/WPA3.}
  \label{fig:4way}
\end{figure}

\subsection{WPA3}

\begin{itemize}
  \item \textbf{\ac{sae}} thay thế cơ chế trao đổi \ac{psk} bằng phương thức
    trao đổi khóa xác thực bằng mật khẩu \cite{wpa3_overview, wpa3spec}. Bản bắt tay đã thu thập không thể bị
    dò ngoại tuyến; mỗi lần thử đều cần tương tác với \ac{ap}, đồng thời khóa
    có tính bảo mật chuyển tiếp.
  \item \textbf{\ac{pmf}} (802.11w) là bắt buộc, giúp chống giả mạo khung
    hủy xác thực và ngắt liên kết.
  \item \textbf{Chế độ WPA3-Enterprise 192 bit} dùng \ac{gcmp}-256 và mật mã
    đường cong elliptic 384 bit cho môi trường yêu cầu bảo mật cao.
  \item \textbf{\ac{owe} (``Enhanced Open'')} \cite{rfc8110, wpa3_deployment} mã hóa mạng mở
    (ví dụ mạng tại quán cà phê) mà không cần mật khẩu, thông qua trao đổi khóa
    không xác thực, qua đó ngăn nghe lén thụ động.
  \item \textbf{\ac{dpp}} (Wi-Fi Easy Connect) hỗ trợ kết nối ban đầu cho
    thiết bị \ac{iot} không có giao diện người dùng bằng mã QR hoặc tương tự.
  \item Các triển khai WPA3 ban đầu có điểm yếu về kênh kề và hạ cấp
    (``Dragonblood''), sau đó được khắc phục bằng các bản cập nhật
    \cite{dragonblood2020, wpa3_deployment}; cần duy trì firmware ở phiên bản mới.
\end{itemize}
Băng tần 6~GHz yêu cầu WPA3 (hoặc \ac{owe}), do đó các mạng Wi-Fi~6E/7 phải
dùng cơ chế bảo mật này.

\subsection{Xác thực trong mạng doanh nghiệp}

IEEE 802.1X \cite{ieee8021x} định nghĩa ba vai trò: \emph{supplicant} (thiết
bị khách cần xác thực), \emph{authenticator} (\ac{ap} hoặc bộ chuyển mạch) và
\emph{máy chủ xác thực}, thường là máy chủ \ac{radius} \cite{rfc2865}. Bộ xác
thực chặn toàn bộ lưu lượng ngoại trừ \ac{eapol} cho đến khi máy chủ xác nhận
thành công; phương thức \ac{eap} \cite{rfc3748} truyền thông tin xác thực bên
trong đường hầm \ac{tls} nếu phương thức đó hỗ trợ.

\begin{table}[htbp]
  \centering\footnotesize
  \caption{Một số phương thức EAP phổ biến.}
  \label{tab:eap}
  \begin{tabularx}{\linewidth}{@{} L{0.6} L{1.2} L{1.2} @{}}
    \toprule
    \textbf{Phương thức} & \textbf{Thông tin xác thực} & \textbf{Ghi chú} \\
    \midrule
    EAP-TLS & Chứng chỉ máy khách và máy chủ & Mạnh nhất; cần hạ tầng khóa công khai (PKI) và phân phối chứng chỉ \\
    \acs{peap} / EAP-TTLS & Chứng chỉ máy chủ; tên người dùng và mật khẩu trong đường hầm & Phổ biến; phải xác minh chứng chỉ máy chủ để tránh lộ thông tin \\
    EAP-FAST & Thông tin xác thực truy cập được bảo vệ & Do Cisco khởi xướng; không yêu cầu chứng chỉ \\
    EAP-SIM / AKA & Thẻ SIM & Chuyển tải Wi-Fi của nhà mạng di động \\
    \bottomrule
  \end{tabularx}
\end{table}

\subsection{Mối đe dọa và biện pháp đối phó}

\begin{table}[!htbp]
  \centering\footnotesize
  \caption{Các mối đe dọa WLAN thường gặp và biện pháp phòng vệ.}
  \label{tab:threats}
  \begin{tabularx}{\linewidth}{@{} L{0.7} L{1.15} L{1.15} @{}}
    \toprule
    \textbf{Mối đe dọa} & \textbf{Mô tả} & \textbf{Biện pháp đối phó} \\
    \midrule
    Nghe lén & Thu thụ động khung trên mạng mở hoặc mã hóa yếu & WPA2/WPA3; \acs{owe} cho mạng khách mở \\
    \acs{ap} trái phép & \acs{ap} không được cấp phép kết nối vào mạng có dây & \acs{wips}, bảo vệ cổng mạng, quét định kỳ \\
    AP giả mạo (evil twin) & \acs{ap} giả bắt chước \acs{ssid} hợp pháp để đánh lừa thiết bị khách & 802.1X có xác minh chứng chỉ máy chủ, \acs{wips} \\
    Tấn công ngắt xác thực/liên kết & Khung quản lý giả khiến thiết bị mất kết nối & \acs{pmf} (802.11w), WPA3 \\
    Dò cụm mật khẩu ngoại tuyến & Dùng từ điển thử với bản bắt tay WPA2 đã thu thập & Cụm mật khẩu dài, ngẫu nhiên; \acs{sae}, 802.1X \\
    Cài đặt lại khóa (\acs{krack}) & Phát lại thông điệp bắt tay để đặt lại nonce \cite{krack2017} & Cập nhật thiết bị khách và \acp{ap} \\
    Lạm dụng mã PIN \acs{wps} & Dò vét cạn mã PIN ngắn của Wi-Fi Protected Setup & Tắt \acs{wps} \\
    Phân đoạn mạng yếu & Thiết bị khách hoặc \ac{iot} bị chiếm quyền truy cập hệ thống nội bộ & \acp{vlan}, tường lửa, cô lập thiết bị khách \\
    \bottomrule
  \end{tabularx}
\end{table}

\subsection{Tăng cường bảo mật}

\begin{itemize}
  \item Dùng WPA3 (hoặc chỉ WPA2-\acs{aes}); tắt \acs{wep},
    \acs{wpa}/\acs{tkip} và \acs{wps}.
  \item Ưu tiên 802.1X dùng chứng chỉ cho mạng doanh nghiệp; với mạng cá nhân/
    \acs{psk}, dùng cụm mật khẩu ngẫu nhiên dài ít nhất 16 ký tự.
  \item Bật \acs{pmf}; cập nhật firmware cho \ac{ap} và bộ điều khiển.
  \item Tách lưu lượng doanh nghiệp, khách, thoại và \ac{iot} vào các \acp{vlan}
    riêng kèm quy tắc tường lửa; bật cô lập thiết bị khách trên \acp{ssid} dành
    cho khách.
  \item Triển khai \acs{wips} hoặc chức năng phát hiện \ac{ap} trái phép; theo
    dõi nhật ký để phát hiện lỗi xác thực lặp lại.
  \item Không dựa vào \acp{ssid} ẩn hoặc lọc \acs{mac}; các biện pháp này gần
    như không tăng bảo mật và dễ dàng bị vượt qua.
\end{itemize}

\begin{tomtat}
\item Ai trong vùng phủ sóng cũng có thể bắt khung, nên \acs{wlan} cần đảm bảo bảo mật, toàn vẹn, xác thực và tính sẵn sàng.
\item Tiến hóa: \acs{wep} (đã bị phá), \acs{wpa}/\acs{tkip}, WPA2 (\acs{aes}-\acs{ccmp}, bắt tay bốn bước; còn rủi ro dò mật khẩu ngoại tuyến và \acs{krack}), WPA3.
\item WPA3: \acs{sae} thay \acs{psk} để chống dò ngoại tuyến, \acs{pmf} bắt buộc, chế độ Enterprise 192-bit, \acs{owe} mã hóa mạng mở, \acs{dpp} đăng ký thiết bị \acs{iot}.
\item Doanh nghiệp dùng 802.1X (supplicant, authenticator, máy chủ \acs{radius}) với \acs{eap}; \acs{eap}-\acs{tls} mạnh nhất, \acs{peap}/\acs{eap}-TTLS phổ biến nhưng phải kiểm tra chứng chỉ máy chủ.
\item Các mối đe dọa chính: nghe lén, rogue \acs{ap}, evil twin, tấn công deauthentication, dò mật khẩu, \acs{krack}, lạm dụng \acs{wps}, phân đoạn yếu.
\item Củng cố: dùng WPA3 (hoặc WPA2-\acs{aes}), tắt \acs{wep}/\acs{tkip}/\acs{wps}, bật \acs{pmf}, dùng mật khẩu ngẫu nhiên dài từ 16 ký tự, phân tách \acs{vlan}, triển khai \acs{wips} và cập nhật firmware. \acs{ssid} ẩn và lọc \acs{mac} không phải là biện pháp bảo mật.
\end{tomtat}
\endgroup
